\chapter{Marchés de prédiction et marchés de paris}\label{ch:m3:prediction-and-betting-markets}

En septembre 2024, un tribunal américain a jugé que le régulateur des produits dérivés avait outrepassé ses pouvoirs en
interdisant à une bourse réglementée d'admettre à la cote des contrats portant sur le parti qui contrôlerait le Congrès. Deux ans
plus tôt, le même régulateur avait infligé une amende à une plateforme \omterm{def:m3:blockchains-for-traders:chain}{blockchain} qui proposait, sans enregistrement, des contrats
comme «~Trump remportera-t-il l'élection présidentielle de 2020~?~». Un contrat qui verse un dollar si un événement se réalise s'échange à un prix qui se lit comme une probabilité.
Les \omterm{def:m3:prediction-and-betting-markets:book}{bookmakers} cotent de tels prix depuis des siècles, avec une marge~; les \omterm{def:m3:prediction-and-betting-markets:exchange}{bourses de paris} permettent désormais aux parieurs de
négocier entre eux, et les régulateurs des produits dérivés ont dû décider lesquels relèvent du jeu d'argent.
Ce chapitre explique comment les cotes encodent probabilités et marges, comment fonctionnent les \omterm{def:m3:prediction-and-betting-markets:exchange}{bourses de paris} et les
\omterm{def:m3:prediction-and-betting-markets:event}{contrats d'événement}, comment les sociétés de trading font ces marchés, et le biais le mieux documenté qui les touche.

\section{Bookmakers, overround et probabilités implicites}

\begin{definition}[Bookmaker, probabilité implicite, overround]\label{def:m3:prediction-and-betting-markets:book}
Un \emph{bookmaker}\index{bookmaker} est une société qui affiche des cotes sur les issues d'un événement et accepte des paris contre elles en qualité de
contrepartie. La \emph{probabilité implicite}\index{probabilité implicite} d'une cote décimale $o$ est $1/o$. L'\emph{overround}\index{overround}
d'un livre est l'écart dont ses probabilités implicites, sommées sur toutes les issues, dépassent un~: la marge du bookmaker.
\end{definition}

Les cotes décimales $o$ versent $o$ par unité misée, mise comprise~; ce sont celles du livre~2, chapitre~29, exprimées
sous forme de montant versé. Un marché à trois issues aux cotes de 1.50, 4.00 et 7.00 présente des \omterm{def:m3:prediction-and-betting-markets:book}{probabilités implicites} de 66.7~\%, 25.0~\% et 14.3~\%, dont la somme vaut 105.95~\%~: un
\omterm{def:m3:prediction-and-betting-markets:book}{overround} de 5.95~\%. Un parieur qui mise sur chaque issue proportionnellement à sa \omterm{def:m3:prediction-and-betting-markets:book}{probabilité implicite} perd 5.6~\% de sa mise quoi qu'il
arrive. Pour lire les probabilités du marché, il faut retirer la marge, et la façon de la retirer relève d'un modèle de l'endroit où le \omterm{def:m3:prediction-and-betting-markets:book}{bookmaker}
l'a répartie.

\begin{proposition}[Quatre façons de retirer la marge]\label{prop:m3:prediction-and-betting-markets:removal}
Soit $\pi_i = 1/o_i$ avec $\sum \pi_i = 1 + m$. La méthode multiplicative pose $p_i = \pi_i/(1 + m)$~; la méthode additive $p_i = \pi_i -
m/n$~; la méthode de puissance $p_i = \pi_i^k$ avec $k > 1$ telle que $\sum p_i = 1$~; la méthode de Shin résout pour la part $z$ des parieurs qui sont
initiés, avec $p_i = \big(\sqrt{z^2 + 4(1 - z)\pi_i^2/(1 + m)} - z\big)/(2(1 - z))$ et $\sum p_i = 1$. La méthode multiplicative répartit la
marge proportionnellement à la probabilité~; les autres en placent relativement davantage sur les outsiders.
\end{proposition}

\begin{proof}
Chaque formule somme à un par construction (les paramètres de puissance et de Shin se calculent par une recherche de racine unidimensionnelle, unique puisque les sommes sont
monotones en ce paramètre). Pour la méthode de puissance, $\pi_i^k/\pi_i = \pi_i^{k - 1}$ est minimal pour le plus petit $\pi_i$, de sorte que les outsiders
perdent le plus de probabilité par rapport à leur valeur implicite~; la formule de Shin fait de même, car le terme des initiés pénalise les petits
$\pi_i$ proportionnellement davantage.
\end{proof}

\begin{table}[ht]
\centering
\small
\begin{tabular}{@{}lrrr@{}}
\toprule
Méthode & favori (1.50) & milieu (4.00) & outsider (7.00) \\
\midrule
Implicite, avec marge & 66.67 & 25.00 & 14.29 \\
Multiplicative & 62.92 & 23.60 & 13.48 \\
Additive & 64.68 & 23.02 & 12.30 \\
Puissance & 64.82 & 22.71 & 12.48 \\
Shin ($z = 3.0\%$) & 64.23 & 23.15 & 12.62 \\
\bottomrule
\end{tabular}
\caption{Probabilités, en pourcentage, retrouvées à partir d'un livre à trois issues avec un \omterm{def:m3:prediction-and-betting-markets:book}{overround} de 5.95~\% par les méthodes de la
\cref{prop:m3:prediction-and-betting-markets:removal}.}\label{tab:m3:prediction-and-betting-markets:methods}
\end{table}

\section{Bourses de paris}

\begin{definition}[Bourse de paris, pari back, pari lay, paris en direct]\label{def:m3:prediction-and-betting-markets:exchange}
Une \emph{bourse de paris}\index{bourse de paris} est une plateforme sur laquelle les parieurs négocient des paris entre eux par l'intermédiaire d'un carnet d'ordres, et l'opérateur perçoit une commission au lieu de prendre des positions.
Un \emph{pari back}\index{pari back} est un pari sur la survenue d'une issue~; un \emph{pari
lay}\index{pari lay} est un pari sur sa non-survenue, celui qui le prend jouant le rôle du \omterm{def:m3:prediction-and-betting-markets:book}{bookmaker} pour ce pari. Les
\emph{paris en direct}\index{paris en direct} sont des paris placés pendant le déroulement de l'événement, à des prix qui évoluent avec lui.
\end{definition}

Une \omterm{def:m3:prediction-and-betting-markets:exchange}{bourse de paris} est un carnet d'ordres à cours limité (livre~1, chapitre~19) dans lequel l'instrument est une issue et le prix est une cote décimale.
Les backers veulent des cotes élevées, les layers des cotes basses, et la fourchette entre la meilleure cote back et la meilleure cote lay est la fourchette du marché. La commission est prélevée sur les gains nets de
chaque marché, de sorte que le parieur ne paie que sur ce qu'il gagne~: Betfair la calcule comme les gains nets multipliés par le taux de base du marché et par un moins la
remise du parieur, les taux de base figurant dans ses pages d'aide australiennes étant de 6~\% pour le sport, 10~\% pour une compétition de rugby à XIII et 8 ou 10~\% pour les
courses locales. Les positions lay et back se compensent~: un parieur qui a misé back à 3.0 puis lay à 2.5 après que les cotes ont baissé verrouille un gain
quoi qu'il arrive, comme un trader qui clôture une position longue à un prix plus élevé. Les marchés en direct évoluent au gré de chaque événement sur le
terrain et récompensent la rapidité de l'information, ce qui place les layers lents à la merci des backers rapides.

\section{Contrats d'événement sur chaîne et hors chaîne}

\begin{definition}[Contrat d'événement, source de résolution]\label{def:m3:prediction-and-betting-markets:event}
Un \emph{contrat d'événement}\index{contrat d'événement} est un produit dérivé qui verse un montant fixe, généralement un dollar, si un événement précisé se réalise, et rien
sinon~; son prix, exprimé en part du montant versé, se lit comme la probabilité que le marché attribue à l'événement. Sa \emph{source de résolution}\index{source de résolution}
est la source et la procédure désignées dans le contrat, par lesquelles la survenue de l'événement est déterminée.
\end{definition}

Aux États-Unis, les \omterm{def:m3:prediction-and-betting-markets:event}{contrats d'événement} sont des dérivés relevant du Commodity Exchange Act et doivent être admis à la cote sur un
\omterm{def:m3:getting-access-commodities-and-power:dcm}{designated contract market} (DCM) (voir le \cref{ch:m3:getting-access-commodities-and-power}). La règle spéciale de la loi permet au régulateur d'interdire les \omterm{def:m3:prediction-and-betting-markets:event}{contrats d'événement} portant sur certaines
activités, parmi lesquelles l'activité illégale et le jeu d'argent, lorsqu'elles sont contraires à l'intérêt public.

\begin{dated}{2026-09}{Contrats d'événement devant les tribunaux et régulateurs américains}\label{dat:m3:prediction-and-betting-markets:kalshi}
Le 12 septembre 2024, le tribunal fédéral de district pour le District de Columbia, dans l'affaire KalshiEX LLC v.\ CFTC, a jugé que les contrats de Kalshi sur le contrôle du Congrès
(«~cette chambre du Congrès sera-t-elle contrôlée par ce parti pour ce mandat~?~») ne relèvent ni de l'activité illégale ni du jeu d'argent, et que l'ordonnance de la CFTC les interdisant
excédait son pouvoir légal, et a fait droit à Kalshi par jugement sommaire. La CFTC a fait appel~; le 7 mai 2025, la cour d'appel a accordé
sa demande non contestée de rejet de l'appel. En janvier 2022, la CFTC avait réglé des accusations contre Blockratize, Inc., exerçant sous le nom de Polymarket, pour
avoir proposé des options binaires sur événements hors bourse sans être un \omterm{def:m3:getting-access-commodities-and-power:dcm}{designated contract market} ni une swap execution facility (SEF) enregistrée~; l'ordonnance imposait
une amende de 1.4~million de dollars et l'obligeait à liquider ses marchés non conformes. Le 2 septembre 2025, les services de la CFTC ont accordé une position de non-action
concernant la déclaration des swaps à QCX, un \omterm{def:m3:getting-access-commodities-and-power:dcm}{designated contract market} et une chambre de compensation que Polymarket avait acquise pour son retour aux États-Unis.
\end{dated}

Sur une plateforme \omterm{def:m3:blockchains-for-traders:chain}{blockchain}, le \omterm{def:m3:prediction-and-betting-markets:event}{contrat d'événement} prend la forme d'une paire de \omterm{def:m3:blockchains-for-traders:key}{jetons}, YES et NO, émis ensemble contre un dollar de \omterm{def:m3:blockchains-for-traders:stable}{stablecoin} et remboursables contre celui-ci,
négociés sur un carnet d'ordres ou auprès d'un teneur de marché, et réglés par un \omterm{def:m3:blockchains-for-traders:contract}{oracle} (voir le \cref{ch:m3:blockchains-for-traders}). La \omterm{def:m3:prediction-and-betting-markets:event}{source de résolution} est la
clause la plus importante et la moins visible du contrat~: «~qui a remporté l'élection~» doit être tranché par quelqu'un, à un moment donné, à partir d'une source, et les règles applicables aux issues contestées
ou ambiguës déterminent ce que versent les \omterm{def:m3:blockchains-for-traders:key}{jetons}. L'émission de paires gonfle aussi les volumes~: à partir du registre on-chain complet de Polymarket pour les
marchés présidentiels de 2024, Tsang et Yang mesurent le volume échangé en octobre sur le marché Trump à 391~millions de dollars, contre 958~millions de volume naïf, qui comptabilise
la création et la destruction de paires comme des transactions.

\begin{dated}{2026-09}{Un oracle optimiste}\label{dat:m3:prediction-and-betting-markets:uma}
La documentation d'UMA décrit son \omterm{def:m3:blockchains-for-traders:contract}{oracle} optimiste~: des proposants tiers déposent une résolution accompagnée d'une caution, remboursable si la proposition est exacte~; la
proposition reste en attente pendant une période de contestation, durant laquelle chacun peut la contester en déposant une caution de contestation~; les litiges sont soumis à un vote des stakers de \omterm{def:m3:blockchains-for-traders:key}{jetons} UMA,
avec une phase de commit de 24 heures et une phase de reveal de 24 heures, et la caution de la partie ayant tort indemnise la partie ayant raison.
\end{dated}

\begin{omfigure}
\centering
\begin{tikzpicture}[font=\small,
  s/.style={draw,thick,rounded corners=2pt,align=center,minimum height=0.9cm,text width=2.5cm},
  arr/.style={-{Stealth[length=2.2mm]},thick}]
\node[s,draw=omDef,fill=omDef!8] (u) at (0,0) {un dollar de\\stablecoin};
\node[s,draw=omThm,fill=omThm!8] (m) at (3.6,0) {émettre un YES\\et un NO};
\node[s,draw=omProp,fill=omProp!8] (t) at (7.2,0) {négocier sur un carnet\\ou un pool};
\node[s,draw=omExo,fill=omExo!8] (o) at (10.8,0) {l'oracle tranche~:\\le gagnant verse \$1};
\draw[arr] (u) -- (m); \draw[arr] (m) -- (t); \draw[arr] (t) -- (o);
\draw[arr,dashed] (m.south) -- ++(0,-0.6) -| (u.south);
\node[font=\footnotesize,anchor=north] at (1.8,-1.1) {un YES et un NO ensemble se remboursent contre un dollar à tout moment};
\end{tikzpicture}

\omcaption{Un \omterm{def:m3:prediction-and-betting-markets:event}{contrat d'événement} sur une plateforme \omterm{def:m3:blockchains-for-traders:chain}{blockchain}~: des \omterm{def:m3:blockchains-for-traders:key}{jetons} YES et NO sont émis ensemble contre un dollar, se négocient séparément, et le \omterm{def:m3:blockchains-for-traders:key}{jeton} gagnant verse
un dollar lorsque l'\omterm{def:m3:blockchains-for-traders:contract}{oracle} tranche l'événement. Comme une paire se rembourse toujours contre un dollar, les cours de YES et de NO totalisent environ un. Schéma.}\label{fig:m3:prediction-and-betting-markets:tokens}
\end{omfigure}

\section{Comment les sociétés de trading animent ces marchés}

Un teneur de marché sur les \omterm{def:m3:prediction-and-betting-markets:event}{contrats d'événement} cote une probabilité avec une fourchette, comme tout autre teneur de marché, et ses risques sont les mêmes sous un habit neuf~: la sélection
adverse exercée par les parieurs qui en savent davantage (livre~1, chapitre~19), un inventaire qu'il ne peut couvrir que sur le même événement, sur d'autres plateformes, et des cours
qui bondissent à zéro ou à un lors de la résolution. Trois traits sont particuliers. Le même événement se négocie sur plusieurs plateformes~: un \omterm{def:m3:prediction-and-betting-markets:book}{bookmaker}, une bourse, une
bourse de \omterm{def:m3:prediction-and-betting-markets:event}{contrats d'événement} réglementée et une plateforme \omterm{def:m3:blockchains-for-traders:chain}{blockchain}, avec des frais, des participants et des \omterm{def:m3:prediction-and-betting-markets:event}{sources de résolution} différents, de sorte que le premier edge d'une société tient à l'\omterm{def:m3:spot-markets:cross}{arbitrage entre plateformes}~:
acheter l'événement là où il est bon marché et le vendre là où il est cher. Le capital est immobilisé jusqu'à la résolution, souvent pendant des mois, si bien qu'un écart de prix doit dépasser le
rendement de ce capital. Et les \omterm{def:m3:prediction-and-betting-markets:event}{sources de résolution} diffèrent~: le même événement phare peut se résoudre différemment sur deux plateformes si leurs sources ou leurs règles pour
les cas limites diffèrent, ce qui transforme un arbitrage en pari sur les règles.

\begin{proposition}[Quand un écart entre plateformes rapporte]\label{prop:m3:prediction-and-betting-markets:gap}
Soit la plateforme A qui vend YES au prix $y_A$ et la plateforme B qui vend NO au prix $1 - y_B$ avec $y_B > y_A$, avec des frais $f_A$ et $f_B$ par contrat et un coût du capital $r$ par an pendant
$\tau$ années jusqu'à la résolution. Acheter un de chaque type rapporte exactement un dollar quoi qu'il arrive, si les deux plateformes résolvent de la même manière, et dégage
$1 - (y_A + 1 - y_B + f_A + f_B)(1 + r\tau)$. Elle est rentable lorsque l'écart $y_B - y_A$ dépasse environ $f_A + f_B + r\tau$.
\end{proposition}

\begin{proof}
Un YES et un NO sur le même événement rapportent ensemble un dollar~; le coût est la somme des deux cours, des frais et du financement jusqu'au paiement.
\end{proof}

\section{Le biais favori-outsider}

\begin{definition}[Biais favori-outsider]\label{def:m3:prediction-and-betting-markets:flb}
Le \emph{biais favori-outsider}\index{biais favori-outsider} est la tendance, documentée sur l'ensemble des marchés de paris, selon laquelle les paris sur les outsiders
rapportent moins par unité misée que les paris sur les favoris, de sorte que les cours des outsiders surestiment leurs probabilités.
\end{definition}

Snowberg et Wolfers, étudiant les courses hippiques, ont cherché à savoir si le biais vient de parieurs qui aiment le risque ou de parieurs qui perçoivent mal les probabilités, et
concluent que les erreurs de perception, comme le prédit la théorie des perspectives, l'expliquent mieux. Le modèle de Shin, dans lequel le \omterm{def:m3:prediction-and-betting-markets:book}{bookmaker} se protège des initiés en décotant davantage
les outsiders, est une autre explication, et la raison d'être de la méthode de Shin présentée dans la \cref{prop:m3:prediction-and-betting-markets:removal}. Pour un trader le
biais est un edge structurel d'un côté~: vendre les outsiders et acheter les favoris, à grande échelle et avec discipline, et une mise en garde~: les cours des outsiders
d'un marché sont de mauvaises estimations de probabilités.

\begin{omfigure}
\begin{tikzpicture}
\begin{axis}[width=11cm,height=5cm,ybar,bar width=14pt,
  xlabel={cote décimale},ylabel={rendement par unité misée, \%},
  tick label style={font=\small},label style={font=\small},
  xmin=0.4,xmax=7.6,xtick={1,2,3,4,5,6,7},xticklabels={1--2,2--3,3--5,5--10,10--20,20--50,50--1000},
  ymin=-45,ymax=0,grid=major,grid style={gray!25},
  table/col sep=comma]
\addplot[fill=omThm!60,draw=omThm] table[x=idx,y=ret]{figdata/markets-3/27-prediction-and-betting-markets/flb.csv};
\end{axis}
\end{tikzpicture}

\omcaption{Rendement réalisé par unité misée selon la tranche de cote, sur un carnet synthétique de 20\,000 courses à huit partants dont les cotes du \omterm{def:m3:prediction-and-betting-markets:book}{bookmaker} donnent des \omterm{def:m3:prediction-and-betting-markets:book}{probabilités
implicites} proportionnelles à la probabilité vraie élevée à la puissance 0.92, avec une marge de 10~\%. Les favoris perdent quelques pour cent~; les outsiders entre 50 et 1\,000 perdent plus de 40~\%.
Résultats synthétiques. Données~: le tutoriel du chapitre.}\label{fig:m3:prediction-and-betting-markets:flb}
\end{omfigure}

\section{Tutoriel~: cotes, marges et biais}\label{tut:m3:prediction-and-betting-markets:odds}

\begin{tutorial}
\textbf{Objectif.} Convertir les cotes des \omterm{def:m3:prediction-and-betting-markets:book}{bookmakers} en probabilités avec les méthodes multiplicative, additive, puissance et de Shin, les comparer, et mesurer un
\omterm{def:m3:prediction-and-betting-markets:flb}{biais favori-outsider} sur un carnet synthétique de résultats. \textbf{État final~:} le \cref{tab:m3:prediction-and-betting-markets:methods},
la \cref{fig:m3:prediction-and-betting-markets:flb} et les chiffres du problème du week-end.
\begin{tutsteps}
\item \textbf{Méthode de Shin.} Une recherche de racine unidimensionnelle pour la proportion d'initiés.
\omcode{code/firm/odds/firm_odds.py}{52}{61}{Suppression de la marge de Shin.}\label{lst:m3:prediction-and-betting-markets:shin}
\item \textbf{Le biais.} Coter des courses avec un \omterm{def:m3:prediction-and-betting-markets:book}{bookmaker} biaisé, tirer les gagnants et moyenner les rendements par cote.
\omcode{code/markets-3/27-prediction-and-betting-markets/python/m3_betting.py}{27}{44}{Un carnet synthétique de courses et les rendements réalisés de ses paris.}\label{lst:m3:prediction-and-betting-markets:flb}
\item \textbf{Exécuter} \texttt{methods()}, \texttt{returns\_by\_odds(flb\_book())}, \texttt{election\_gap()} et \texttt{fig\_betting.py}.
\end{tutsteps}
\textbf{Pour aller plus loin.} Retrouver le paramètre de biais à partir des résultats synthétiques par maximum de vraisemblance~; couvrir un \omterm{def:m3:prediction-and-betting-markets:exchange}{pari back} vers un résultat égal avec
\texttt{hedge\_equal}~; ajouter une trajectoire de cotes en direct et calculer le risque d'inventaire d'un teneur de marché.
\end{tutorial}

\section{Construction~: cotes et couverture}\label{bld:m3:prediction-and-betting-markets:odds}

\begin{build}
\bfield{Rôle} La société miniature cote les événements à travers les \omterm{def:m3:prediction-and-betting-markets:book}{bookmakers}, les bourses et les plateformes de \omterm{def:m3:prediction-and-betting-markets:event}{contrats d'événement}~: elle a besoin de leurs cours sous forme de probabilités sans
marges, de frais exprimés sur la même base, et de couvertures qui égalisent ses résultats.
\bfield{Interface} \texttt{implied}, \texttt{overround}~; \texttt{multiplicative}, \texttt{additive}, \texttt{power}, \texttt{shin}~; \texttt{kelly(p, odds)}~;
\texttt{exchange\_payout(stake, odds, commission)}~; \texttt{hedge\_equal(back\_stake, back\_odds, lay\_odds, commission)}.
\bfield{Règles} Cotes décimales en interne, conversions aux interfaces~; probabilités dont la somme vaut un~; commission appliquée sur les gains nets, comme le font les bourses.
\bfield{Tests d'acceptation} \texttt{code/firm/odds/tests/}~: l'\omterm{def:m3:prediction-and-betting-markets:book}{overround}~; toutes les méthodes donnent une somme égale à un~; les méthodes puissance et de Shin décotent davantage les outsiders que la
méthode multiplicative~; Kelly à edge nul et à un edge donné~; une couverture vers des résultats égaux.
\bfield{Pour aller plus loin} Détecteur d'arbitrage multi-plateformes pour un même événement~; la décote pour risque de résolution entre plateformes aux sources différentes.
\end{build}

\begin{omsources}
\item KalshiEX LLC v.\ CFTC, n°~23-3257 (D.D.C.), note motivée du 12 septembre 2024 (Doc 51)~; D.C. Circuit n°~24-5205, ordonnance du 7 mai 2025
accueillant la demande de rejet non contestée de la CFTC.
\item CFTC, communiqué de presse 8478-22 (Polymarket), janvier 2022~; CFTC Letter n°~25-28, 2 septembre 2025~; CoinDesk, 3 septembre 2025.
\item Betfair, pages d'aide sur la commission («~Exchange: What is Commission and how is it calculated?~»~; version australienne «~Commissions and Charges~»),
copies de l'Internet Archive.
\item K.~P.~Tsang et Z.~Yang, «~The Anatomy of a Blockchain Prediction Market: Polymarket in the 2024 U.S. Presidential Election~»,
arXiv:2603.03136, version du 14 septembre 2026.
\item Documentation UMA, «~How does UMA's oracle work~». Consulté en septembre 2026.
\item H.~S.~Shin, «~Measuring the incidence of insider trading in a market for state-contingent claims~», \emph{Economic Journal} 103, 1993~; E.~Snowberg
et J.~Wolfers, «~Explaining the Favorite-Longshot Bias: Is it Risk-Love or Misperceptions?~», NBER Working Paper 15923, 2010.
\end{omsources}

\section{Exercices}

\begin{exercise}[$\star$]\label{exo:m3:prediction-and-betting-markets:1}
Un marché à deux issues est coté à 1.80 et 2.10. Calculer son \omterm{def:m3:prediction-and-betting-markets:book}{overround} et les probabilités multiplicatives.
\end{exercise}

\begin{exercise}[$\star$]\label{exo:m3:prediction-and-betting-markets:2}
Un parieur a pris un \omterm{def:m3:prediction-and-betting-markets:exchange}{pari back} à 3.0 pour 100, et la cote est passée à 2.5. Calculer la mise à engager en \omterm{def:m3:prediction-and-betting-markets:exchange}{pari lay} pour garantir un gain égal, et ce gain.
\end{exercise}

\begin{exercise}[$\star$]\label{exo:m3:prediction-and-betting-markets:3}
Expliquer pourquoi une bourse américaine qui admet à la cote des \omterm{def:m3:prediction-and-betting-markets:event}{contrats d'événement} doit être un \omterm{def:m3:getting-access-commodities-and-power:dcm}{designated contract market} (DCM).
\end{exercise}

\begin{exercise}[$\star\star$]\label{exo:m3:prediction-and-betting-markets:4}
Calculer les probabilités de la méthode de la puissance pour les cotes 1.50, 4.00 et 7.00, et expliquer pourquoi celle de l'outsider est inférieure à la probabilité multiplicative.
\end{exercise}

\begin{exercise}[$\star\star$]\label{exo:m3:prediction-and-betting-markets:5}
Un \omterm{def:m3:prediction-and-betting-markets:event}{contrat d'événement} se négocie à 0.58 et vous estimez la probabilité à 0.65. Quelle fraction de Kelly faut-il miser, sans tenir compte des frais~?
\end{exercise}

\begin{exercise}[$\star\star$]\label{exo:m3:prediction-and-betting-markets:6}
Pourquoi un arbitrage entre les contrats de deux plateformes portant sur «~le même~» événement comporte-t-il un risque de résolution~?
\end{exercise}

\begin{exercise}[$\star\star\star$]\label{exo:m3:prediction-and-betting-markets:7}
\emph{Programmation.} Avec \texttt{returns\_by\_odds}, quel est le rendement d'un \omterm{def:m3:prediction-and-betting-markets:exchange}{pari back} sur chaque concurrent dont la cote est comprise entre 1 et 3, et sur chaque concurrent dont la cote est comprise entre
50 et 1\,000~?
\end{exercise}

\begin{exercise}[$\star\star\star$]\label{exo:m3:prediction-and-betting-markets:8}
\emph{Repérer l'erreur.} «~Le marché cote l'outsider à 50 contre 1, donc sa probabilité vaut 2~\%.~»
\end{exercise}

\section{Problème~: Soirée électorale}

\begin{problem}[{Problème du week-end --- un événement, deux plateformes}]\label{pb:m3:prediction-and-betting-markets:1}
Un mois avant une élection, la plateforme A (une bourse de \omterm{def:m3:prediction-and-betting-markets:event}{contrats d'événement} régulée) vend le YES sur un candidat à 0.58 avec des frais d'un cent par contrat~; la plateforme B (une
plateforme \omterm{def:m3:blockchains-for-traders:chain}{blockchain}) vend le YES à 0.62, donc le NO à 0.38, avec des frais d'un demi-cent. Le capital rapporte ailleurs 4.5~\% par an. Le modèle attribue au candidat 65~\%.

\medskip\noindent\textbf{Partie I --- L'arbitrage.}
\begin{enumerate}
\item Quel est le coût d'un YES sur A et d'un NO sur B, frais compris~?
\item Que rapporte la paire, et quand~?
\item Quel est le profit par paire après le coût du capital~?
\item Quel est le plus petit écart qui résiste aux frais et au coût du capital~?
\item Que rapportent 10\,000 paires, et quel capital immobilisent-elles~?
\end{enumerate}

\medskip\noindent\textbf{Partie II --- Les risques.}
\begin{enumerate}[resume]
\item Que se passe-t-il si les deux plateformes ne résolvent pas l'événement de la même façon~?
\item Qu'ajoute le règlement-livraison de la plateforme B en \omterm{def:m3:blockchains-for-traders:stable}{stablecoins}~?
\item Qu'est-ce qui limite la taille de la transaction sur chaque plateforme~?
\item Une société résidant aux États-Unis peut-elle négocier sur les deux plateformes~?
\item Que révèle l'écart sur les participants des deux plateformes~?
\end{enumerate}

\medskip\noindent\textbf{Partie III --- L'opinion.}
\begin{enumerate}[resume]
\item Avec une conviction de 65~\%, quelle fraction de Kelly faudrait-il miser sur le YES de la plateforme A~?
\item Quel est le rendement espéré par contrat à 0.58~?
\item Pourquoi un professionnel miserait-il une fraction de Kelly~?
\item Comment vérifier si 65~\% est meilleur que les 58 à 62~\% du marché~?
\item Quelle incidence le \omterm{def:m3:prediction-and-betting-markets:flb}{biais favori-outsider} a-t-il sur l'opinion~?
\end{enumerate}

\medskip\noindent\textbf{Partie IV --- Jugement.}
\begin{enumerate}[resume]
\item Les contrats sur élections sont-ils des instruments de jeu ou de couverture~?
\item Qu'a tranché la décision de justice de 2024, et que laisse-t-elle en suspens~?
\item Le cours de quelle plateforme est la meilleure prévision, et comment le savoir~?
\item Énoncer le \emph{résultat nommé}~: l'écart de prix qui résiste aux frais et au coût du capital, et la mise de Kelly sur le contrat le moins cher pour une conviction de 65~\%.
\item En une phrase~: que représente un cours de 0.58~?
\end{enumerate}
\end{problem}

\section{Questions d'entretien}

\begin{interviewq}[$\star$ \iqroles{trader}]\label{iq:m3:prediction-and-betting-markets:1}
Comment convertiriez-vous la cote d'un \omterm{def:m3:prediction-and-betting-markets:book}{bookmaker} en probabilités~?
\end{interviewq}

\begin{interviewq}[$\star$ \iqroles{chercheur}]\label{iq:m3:prediction-and-betting-markets:2}
Qu'est-ce que le \omterm{def:m3:prediction-and-betting-markets:flb}{biais favori-outsider} et qu'est-ce qui l'explique~?
\end{interviewq}

\begin{interviewq}[$\star\star$ \iqroles{trader}]\label{iq:m3:prediction-and-betting-markets:3}
Comment assureriez-vous la teneur de marché d'un \omterm{def:m3:prediction-and-betting-markets:event}{contrat d'événement} un mois avant l'événement~?
\end{interviewq}

\begin{interviewq}[$\star\star$ \iqroles{risque}]\label{iq:m3:prediction-and-betting-markets:4}
Quels sont les risques liés à la détention de positions opposées sur un même événement, sur deux plateformes~?
\end{interviewq}

\begin{interviewq}[$\star\star$ \iqroles{chercheur}]\label{iq:m3:prediction-and-betting-markets:5}
Établissez la mise de Kelly pour un pari à cote décimale $o$ et probabilité de gain $p$.
\end{interviewq}

\begin{interviewq}[$\star\star\star$ \iqroles{développeur, risque}]\label{iq:m3:prediction-and-betting-markets:6}
Concevez un \omterm{def:m3:blockchains-for-traders:contract}{oracle} permettant de résoudre les \omterm{def:m3:prediction-and-betting-markets:event}{contrats d'événement}, difficile à manipuler.
\end{interviewq}
